\chapter{Entering the advanced album: algebra, quantum states, and fluids}
\label{learn:advanced}

\begin{learningcard}{Prerequisites and three destinations}
Use vector spaces and inner products, probability, algebraic structure, and
numerical conservation from the preceding chapters. We now construct one
bounded entry example for each of three advanced routes: Cayley--Dickson
algebras, quantum states, and lattice-Boltzmann computation. The examples
supply enough vocabulary and checks to read the advanced review critically;
they leave specialized classification, quantum information, and fluid
analysis as further coursework.
\end{learningcard}

\section{Doubling changes multiplication as well as dimension}
\label{learn:hypercomplex}
Starting with $\RR$ and its identity conjugation, the Cayley--Dickson
construction replaces an algebra by ordered pairs. Fix the convention
\[
 (a,b)(c,d)=(ac-\overline d\,b,\;da+b\overline c),
 \qquad \overline{(a,b)}=(\overline a,-b).
\]
The real vector-space dimension doubles at each step. Basis vectors are
ordered recursively: the old basis occupies the first half, followed by
pairs with zero first component and old basis in the second half. With
$e_0=1$, this fixes every multiplication sign used below.
Successive dimensions one, two, four, and eight give
$\RR,\CC,\HH,\OO$. The product loses commutativity at quaternions and
associativity at octonions. Norm multiplicativity survives through octonions
but fails for sedenions in dimension sixteen \citep{baez2001octonions}.

\section{Worked example: quaternion arithmetic and a boundary}
Write a quaternion as $q=a+bi+cj+dk$, with
$i^2=j^2=k^2=-1$, $ij=k$, $jk=i$, $ki=j$, and reversed distinct products
having the opposite sign. Then
\[
 (1+i)(1+j)=1+i+j+k,\qquad
 (1+j)(1+i)=1+i+j-k.
\]
Conjugation sends $q$ to $\overline q=a-bi-cj-dk$, and
$q\overline q=a^2+b^2+c^2+d^2=\|q\|^2$. For $q\ne0$,
$q^{-1}=\overline q/\|q\|^2$. Associativity and
$\overline{pq}=\overline q\,\overline p$ give
\[
 \|pq\|^2=pq\overline q\,\overline p
 =\|q\|^2p\overline p=\|p\|^2\|q\|^2.
\]
The scalar interpretation of the products justifies the norm equation.
For the displayed example, each factor has norm $\sqrt2$ and the product
has norm two. Noncommutativity coexists with a multiplicative norm.

A zero divisor is a nonzero element whose product with another nonzero
element is zero. Under the fixed doubling convention, the sedenion vectors
$a=e_3+e_{10}$ and $b=e_6-e_{15}$ satisfy
\[
 e_3e_6=e_5,\quad e_3e_{15}=e_{12},\quad
 e_{10}e_6=e_{12},\quad e_{10}e_{15}=e_5,
 \qquad ab=e_5-e_{12}+e_{12}-e_5=0.
\]
Each factor has Euclidean norm $\sqrt2$, demonstrating norm failure as well
as zero divisors. The coordinate convention is essential to replaying this
witness. It also explains why extending a successful quaternion construction
to every doubled algebra requires fresh proofs of the properties used.

\begin{learningcard}{Historical situation: reconstruct a mathematical route}
Baez's survey of the octonions describes the progression through normed
division algebras and their historical development \citep{baez2001octonions}.
It is a modern scholarly account rather than a contemporaneous record of
an inventor's thoughts. Our doubling sequence is an instructional
reconstruction: it makes the algebraic dependencies visible without
asserting that every historical participant followed this exact path.
\end{learningcard}

\section{Worked example: probabilities from a complex state}
\label{learn:quantum}
A pure qubit is represented by a unit vector
$\psi=\alpha|0\rangle+\beta|1\rangle$ in $\CC^2$, with
$|0\rangle=(1,0)^T$, $|1\rangle=(0,1)^T$ and
$|\alpha|^2+|\beta|^2=1$. Vectors differing by a common phase
$e^{i\theta}$ represent the same pure state. The complex inner product is
$\langle\phi,\psi\rangle=\sum_j\overline{\phi_j}\psi_j$.
The Born rule assigns computational-basis outcome probabilities
$P(0)=|\alpha|^2$ and $P(1)=|\beta|^2$.

A unitary matrix $U$ satisfies $U^\dagger U=I$, where the dagger denotes
conjugate transpose. It preserves inner products and normalization because
$\langle U\phi,U\psi\rangle=\phi^\dagger U^\dagger U\psi
=\langle\phi,\psi\rangle$. For the Hadamard matrix
\[
 H=\frac1{\sqrt2}\begin{pmatrix}1&1\\1&-1\end{pmatrix},\qquad
 H|0\rangle=\frac{|0\rangle+|1\rangle}{\sqrt2},
\]
the two measurement probabilities are $1/2$. Applying $H$ again gives
$H^2|0\rangle=|0\rangle$, so the final probability of zero is one.
Amplitudes combine before their magnitudes are squared. Replacing the
intermediate state by a classical random bit would erase relative-phase
information relevant to later operations.

The state-space postulates and Born rule are physical modeling rules, while
normalization preservation is a mathematical consequence of unitarity.
Connecting either to an experiment requires a specified preparation, operation,
and measurement. A classical program holding complex arrays does not by
itself demonstrate a quantum device. For the formal introductory framework,
see Nielsen and Chuang, Chapter 2.\footnote{Michael A. Nielsen and Isaac L.
Chuang, \emph{Quantum Computation and Quantum Information}, 10th anniversary
edition (2010), Chapter 2, especially the postulates of quantum mechanics:
\url{https://doi.org/10.1017/CBO9780511976667}.}

\section{A lattice-Boltzmann step has exact moment obligations}
\label{learn:fluids}
A lattice-Boltzmann method evolves populations $f_i$ associated with a finite
set of velocities $c_i$. In lattice units, define density and momentum by
$\rho=\sum_i f_i$ and $\rho u=\sum_i c_i f_i$ for $\rho>0$.
For the square D2Q9 stencil, velocities are $(0,0)$, the four axial unit
vectors, and the four diagonal vectors $(\pm1,\pm1)$. Their weights are
$4/9$, $1/9$ per axial vector, and $1/36$ per diagonal vector.
The weights sum to one; symmetry gives
\[
 \sum_iw_ic_i=0,\qquad
 \sum_iw_i c_{i\alpha}c_{i\beta}=\frac13\delta_{\alpha\beta},
\]
where $\delta_{\alpha\beta}$ is one for equal indices and zero otherwise.
Odd third moments vanish by opposite-velocity pairing.

The common low-velocity equilibrium is
\[
 f_i^{\mathrm{eq}}=w_i\rho\left[
 1+3c_i\cdot u+\frac92(c_i\cdot u)^2-\frac32|u|^2\right].
\]
Summing the polynomial with the weight identities gives
$\sum_i f_i^{\mathrm{eq}}=\rho[1+(9/2)(|u|^2/3)-(3/2)|u|^2]=\rho$.
Multiplying by $c_i$ and summing gives $\sum_i c_if_i^{\mathrm{eq}}=\rho u$:
the linear term contributes $3\rho u/3$, and the other vector sums vanish.
Thus the collision step
$f_i^*=f_i-\omega(f_i-f_i^{\mathrm{eq}})$ preserves both moments in exact
arithmetic. Streaming $f_i(x+c_i,t+1)=f_i^*(x,t)$ permutes populations on a
periodic lattice, so global mass is preserved as well.

Deriving incompressible Navier--Stokes behavior requires further assumptions:
low Mach number, suitable scale separation, the appropriate higher stencil
moments, and compatible boundary and forcing rules. In the standard BGK
hydrodynamic limit, lattice-unit kinematic viscosity is
$\nu=(1/3)(1/\omega-1/2)$; $0<\omega<2$ makes that coefficient positive,
but is insufficient alone to guarantee every simulation's stability.
Diffusion equations can use related scalar-population constructions, with
different equilibrium and moment targets. Naming both approaches
``lattice-Boltzmann'' leaves their modeled quantities distinct.
For a derivation and its assumptions, consult the method review by Chen and
Doolen.\footnote{Shiyi Chen and Gary D. Doolen, ``Lattice Boltzmann Method
for Fluid Flows,'' \emph{Annual Review of Fluid Mechanics} 30 (1998),
pp.~329--364: \url{https://doi.org/10.1146/annurev.fluid.30.1.329}.}

\section{Exercises with full solutions}
\paragraph{1. Invert a quaternion.}
Find the inverse of $q=1+2i$ and check the product.

\paragraph{Solution.}
The squared norm is five, so $q^{-1}=(1-2i)/5$.
Multiplication gives $(1+2i)(1-2i)/5=(1-4i^2)/5=1$.
The reversed product also equals one. This calculation uses the quaternion
relations and proves an inverse for the specified element.

\paragraph{2. Keep relative phase until the measurement.}
Apply $H$ to $(|0\rangle-|1\rangle)/\sqrt2$ and find the outcome probabilities.

\paragraph{Solution.}
Matrix multiplication gives
$\frac12(1-1,1+1)^T=(0,1)^T=|1\rangle$.
The probabilities are zero for outcome zero and one for outcome one.
Before the transformation, both probabilities were $1/2$. The relative
minus sign changes the subsequent interference.

\paragraph{3. Detect a population defect.}
Suppose a faulty equilibrium sums to $\rho+\varepsilon$ at one lattice
site. What does collision do to that site's density?

\paragraph{Solution.}
Summing the update gives
$\sum_i f_i^*=\rho-\omega[\rho-(\rho+\varepsilon)]
=\rho+\omega\varepsilon$.
The defect injects density in proportion to the relaxation factor. Periodic
streaming redistributes that injected amount and preserves its global sum;
it cannot repair the collision defect. This is why moment identities should
be checked before interpreting pictures of a flow.

\section{Read the existing advanced review by its evidence contract}
The advanced mathematical baseline studies exceptional Lie algebras,
Cayley--Dickson properties, modular forms, and fractal dimensions. The
computational portfolio then asks what retained artifacts establish about
implementations and proposed physical mechanisms. Carry three questions into
each entry: which mathematical statement has a proof, which finite calculation
has been reproduced, and which physical inference needs an independently
specified model and observation? Subsequent courses can deepen representation
theory, complex analysis, functional analysis, quantum information, and
continuum mechanics along separate paths. These chapters establish a shared
entry language and worked checks, rather than completion of those subjects.
